% !PRINTED 357
\weberpage{357}
%<<CONT
relativ prim sind, $x$ nicht durch $p$ teilbar sein, und es muß $\mathfrak{a}$
durch $p$ teilbar sein.

Setzen wir $\mathfrak{a} = p^k \mathfrak{a}_1$, wo $\mathfrak{a}_1$ nicht mehr durch $p$ teilbar ist,
und bezeichnen mit $\mathfrak{p}$ ein in $p$ aufgehendes Primideal, so muß
wegen (11) $\mathfrak{p}$ in $\mathfrak{m}$ oder in $\mathfrak{m}'$ aufgehen, und da $\mathfrak{m}$ primär vorausgesetzt war, kann $\mathfrak{p}$ nicht $= p$ sein. Es ist daher $N(\mathfrak{p}) = p = \mathfrak{p}\mathfrak{p}'$,
und wenn $\mathfrak{m}$ durch $\mathfrak{p}$ teilbar ist, kann es nicht zugleich durch $\mathfrak{p}'$
teilbar sein. Folglich ist $\mathfrak{m}$ durch $\mathfrak{p}^k$ teilbar. Es ist also nach
der Definition von $\mathfrak{m}$ (Satz 6) auch $\dfrac{b+\sqrt D}{2}$ durch $\mathfrak{p}^k$ teilbar,
folglich $\dfrac{b+\sqrt D}{2} y$ mindestens durch $\mathfrak{p}^{k+1}$, und $ax$ und folglich $\mu$
genau durch $\mathfrak{p}^k$, während andererseits $\mathfrak{a}\mathfrak{m} = \mu$ gleichfalls mindestens durch $\mathfrak{p}^{k+1}$ teilbar ist, worin ein Widerspruch liegt. Ähnliches gilt auch, wenn $p = \mathfrak{p}^2$ ist, also in $\Delta$ (aber nicht in $Q$)
aufgeht. Folglich kann in $\mathfrak{a}$ keine rationale Primzahl enthalten
sein und $\mathfrak{a}$ ist primär.

Nach dem, was bisher bewiesen ist, können wir noch folgenden
Satz formulieren:

9. \gesperrt{Satz. Sind}
\[
\mu = ax + \frac{b+\sqrt D}{2} y = \mathfrak{m}\mathfrak{a},
\]
\[
\mu' = ax' + \frac{b+\sqrt D}{2} y' = \mathfrak{m}\mathfrak{a}'
\]
% Source scan supports b before \sqrt D, matching the parallel line above.
\gesperrt{zwei zu $Q$ teilerfremde Zahlen, haben weder $x$ und $y$
noch $x'$ und $y'$ einen gemeinschaftlichen Teiler, ist}
\[
N(\mathfrak{a}) = N(\mathfrak{a}') = A,
\]
\gesperrt{und für eine und dieselbe Zahl $B$}
\[
\begin{aligned}
\frac{B+\sqrt D}{2}&\equiv 0 \pmod{\mathfrak{a}}\\
&\equiv 0 \pmod{\mathfrak{a}'},
\end{aligned}
\]

\begin{center}\rule{0.2\linewidth}{0.4pt}\end{center}
\gesperrt{so sind die beiden Ideale $\mathfrak{a}$ und $\mathfrak{a}'$ identisch, und die
Zahlen $\mu$, $\mu'$ unterscheiden sich nur durch einen Einheitsfaktor.}

Denn $\mathfrak{a}$ sowohl als $\mathfrak{a}'$ sind nach 8. primär und sind nach 6.
definiert als größter gemeinschaftlicher Teiler von
\[
A \text{ und } \frac{B+\sqrt D}{2}.
\]
